\documentclass[10pt]{beamer}

% Tema elegante y académico
\usetheme{Madrid}
\usecolortheme{default}

% Paquetes útiles
\usepackage[spanish]{babel}
\usepackage[utf8]{inputenc}
\usepackage{graphicx}
\usepackage{amsmath}

% Información de la portada
\title[Grupo 2]{Proyecto Final: Titulo en construccion}
\author{Jeronimo Gil - Isaac - Oscar Díaz}
\institute{Universidad EAFIT}
\date{\today}

\begin{document}

%------------------------------------------------
% Portada
%------------------------------------------------
\begin{frame}
\titlepage
\end{frame}

%------------------------------------------------
% Índice
%------------------------------------------------
\begin{frame}{Contenido}
\tableofcontents
\end{frame}


%------------------------------------------------
\section{Primer Ejercicio}
%------------------------------------------------

\subsection{Descripcion del Ejercicio}
\begin{frame}{Descripcion del Ejercicio}

Ejercicio 5:

La empresa relojera Kronos invirtió hace 3 años 2.000.000 de francos suizos (CHF) en una subsidiaria en Brasil. La tasa obtenida de rendimiento de la inversión en reales brasileros (BRL) es del 16\(\%\) ASV (Anual Semestre Vencido). Calcule la rentabilidad anual efectiva en francos suizos y la cantidad de francos suizos obtenidos al final del tercer año.

Al momento de la inversion, las tasas de cambio eran las siguientes:
\begin{itemize}
    \item TC CHF/USD = 0.9687 

    \item TC BRL/USD = 4.06275
\end{itemize}

Hoy las tasas de cambio son:

\begin{itemize}
    \item TC CHF/USD = 0.92447 

    \item TC BRL/USD = 5.32004
\end{itemize}

\end{frame}

\subsection{Resultados}


\subsection{Resultados}

%------------------------------------------------
\begin{frame}{Inversión Internacional Kronos (Parte 1)}

\small

\begin{columns}[T,onlytextwidth]

% COLUMNA IZQUIERDA
\begin{column}{0.48\textwidth}

\textbf{Datos del problema}

\begin{itemize}
    \item Inversión inicial: \(C_0=2\,000\,000\) CHF.
    \item Plazo: \(n=3\) años.
    \item Rentabilidad en Brasil: \(i_{ASV}=16\%\).
\end{itemize}


\begin{table}
\centering
\footnotesize
\begin{tabular}{lcc}
\hline
\textbf{Tasa} & \textbf{Inicial} & \textbf{Final}\\
\hline
CHF/USD & 0.9687 & 0.92447\\
BRL/USD & 4.06275 & 5.32004\\
\hline
\end{tabular}
\end{table}

\textbf{1. Tasa efectiva anual en BRL}

Como la tasa es anual semestre vencido:

\[
i_{BRL}=
\left(1+\frac{0.16}{2}\right)^2-1
\]

\[
\boxed{i_{BRL}=16.64\%}
\]

\end{column}

% COLUMNA DERECHA
\begin{column}{0.48\textwidth}

\textbf{2. Tasa de cambio CHF/BRL}

Se divide la tasa CHF/USD entre la tasa BRL/USD:

\[
TC_0=\frac{0.9687}{4.06275}
\]

\[
TC_0\approx 0.2384
\]

\vspace{-0.2cm}

\[
TC_3=\frac{0.92447}{5.32004}
\]

\[
TC_3\approx 0.1738
\]

\end{column}

\end{columns}

\end{frame}

%------------------------------------------------
\begin{frame}{Inversión Internacional Kronos (Parte 2)}

\small

\begin{columns}[T,onlytextwidth]

% COLUMNA IZQUIERDA
\begin{column}{0.48\textwidth}

\textbf{3. Variación de la tasa de cambio}

La variación acumulada durante los tres años es:

\[
\operatorname{Var}_{CHF/BRL}
=\frac{TC_3}{TC_0}-1
\approx -27.10\%
\]

La variación anual efectiva es:

\[
i_{TC}=
(1+\operatorname{Var}_{CHF/BRL})^{1/3}-1
\approx -10.01\%
\]


\end{column}

% COLUMNA DERECHA
\begin{column}{0.48\textwidth}

\textbf{4. Rentabilidad efectiva anual en CHF}

Se combina la rentabilidad en reales con la variación de la tasa de cambio:
 
\[
\boxed{i_{CHF}=(1+i_{BRL})(1+i_{TC})-1 \approx 4.97\%}
\]

\textbf{5. Capitalización de la inversión}

El capital final después de tres años es:

\[
\boxed{C_3=C_0(1+i_{CHF})^3\approx 2\,313\,029.93\ \text{CHF}}
\]

Rentabilidad efectiva anual:

\[
\boxed{i_{CHF}\approx 4.97\%}
\]


\end{column}

\end{columns}

\end{frame}

\subsection{Discusion}
\begin{frame}{Decrecimiento de la tasa de cambio CHF/BRL}
\centering
\includegraphics[width=0.92\textwidth]{images/plot1_tc_chf_brl.png}
\end{frame}

\begin{frame}{Capitalización BRL vs CHF}
\centering
\includegraphics[width=0.92\textwidth]{images/plot2_capitalizacion.png}
\end{frame}

%------------------------------------------------
\section{Segundo Ejercicio}
%------------------------------------------------

\subsection{Descripcion del Ejercicio}
\begin{frame}{Descripcion del Ejercicio}

Ejercicio 6:

La empresa mexicana ”Tumex” esta comprando una mina de caliza en Peru por valor de USD (dolares) 200 Millones. La empresa cuenta con USD 30 millones y el resto lo debe financiar con el sistema financiero a un plazo de 6 anos. Tiene 3 alternativas:

\begin{itemize}
    \item El banco INCA le presta al 10\(\%\) en PEN (nuevos soles peruanos).
    \item El banco AZTECA le presta en MXN (nuevos pesos mexicanos) al 12\(\%\).
    \item El banco APACHE le presta al 6\(\%\) en USD.
\end{itemize}

Las tasas de cambio al momento del desembolso son:
\begin{itemize}
    \item PEN/USD = 3.50
    \item MXN/USD = 22.30
\end{itemize}

La variacion esperada del tipo de cambio es:

\begin{itemize}
    \item PEN/USD = 5.5\(\%\) anual.
    \item MXN/USD = 7\(\%\) anual.
\end{itemize}

\textit{¿Cual es la alternativa mas barata y cuanto es su costo financiero en la moneda de la empresa, que es el MXN?}

\end{frame}


\subsection{Resultados}

%------------------------------------------------
\begin{frame}{Financiamiento Internacional (Parte 1)}

\small

\begin{columns}[T,onlytextwidth]

\begin{column}{0.48\textwidth}

\textbf{1. Datos del problema}

\begin{itemize}
    \item Valor de la mina: USD 200 millones.
    \item Recursos propios: USD 30 millones.
    \item Deuda: USD 170 millones a 6 años.
\end{itemize}

\textbf{Tasas de cambio iniciales}

\[
TC_{\text{PEN/USD}}=3.50,\qquad
TC_{\text{MXN/USD}}=22.30
\]

\textbf{2. Tasa de cambio cruzada}

\[
TC_{\text{MXN/PEN},0}=\frac{22.30}{3.50}
\]

\[
\boxed{TC_{\text{MXN/PEN},0}=6.3714}
\]

\end{column}

\begin{column}{0.48\textwidth}

\textbf{3. Tasas de cambio al año 1}

Con variaciones esperadas del \(5.5\%\) en PEN/USD y \(7\%\) en MXN/USD:

\[
TC_{\text{PEN/USD},1}=3.50(1.055)=3.6925
\]

\[
TC_{\text{MXN/USD},1}=22.30(1.07)=23.861
\]

Por tanto:

\[
TC_{\text{MXN/PEN},1}=\frac{23.861}{3.6925}
\]

\[
\operatorname{Var}_{\text{MXN/PEN}}
=\frac{TC_1}{TC_0}-1
\]

\[
\boxed{\operatorname{Var}_{\text{MXN/PEN}}=1.42\%}
\]

\end{column}

\end{columns}

\end{frame}

%------------------------------------------------
\begin{frame}{Financiamiento Internacional (Parte 2)}

\small

\begin{columns}[T,onlytextwidth]

\begin{column}{0.48\textwidth}

\textbf{4. Costo financiero en MXN}

Para el banco INCA, la tasa del \(10\%\) en PEN se ajusta por la variación cambiaria:

\[
i_{\text{INCA}}=(1+0.10)(1+0.0142)-1
\]

\[
\boxed{i_{\text{INCA}}=11.56\%}
\]

El banco AZTECA presta directamente en MXN al \(12\%\) E.A.:

\[
\boxed{i_{\text{AZTECA}}=12\%}
\]

\end{column}

\begin{column}{0.48\textwidth}

Para el banco APACHE, la tasa del \(6\%\) en USD se ajusta por la variación del \(7\%\) del MXN/USD:

\[
i_{\text{APACHE}}=(1+0.06)(1+0.07)-1
\]

\[
\boxed{i_{\text{APACHE}}=13.42\%}
\]

\textbf{5. Comparación y decisión}

\begin{table}
\centering
\footnotesize
\begin{tabular}{lc}
\hline
\textbf{Banco} & \textbf{Costo E.A. en MXN}\\
\hline
INCA & \(11.56\%\)\\
AZTECA & \(12.00\%\)\\
APACHE & \(13.42\%\)\\
\hline
\end{tabular}
\end{table}

\[
\boxed{\text{Alternativa más barata: INCA}}
\]

\end{column}

\end{columns}

\end{frame}

%------------------------------------------------
\section{Tercer Ejercicio}
%------------------------------------------------

\subsection{Descripcion del Ejercicio}

\begin{frame}{Descripcion del Ejercicio}

\small

\textbf{Objetivo:} Diseñar una estrategia de inversión que permita cubrir las cuotas mensuales y los abonos extraordinarios de un crédito en UVR, utilizando un capital inicial de \$760.000.000 COP.

\textbf{Condiciones del crédito}

\begin{itemize}
    \item \textbf{Monto:} 1.000.000 UVR.
    \item \textbf{Plazo:} 120 meses (10 años).
    \item \textbf{Tasa de interés:} UVR + 10\% E.A.
    \item \textbf{Amortización:} Cuotas fijas mensuales en UVR.
    \item \textbf{Período de gracia:} 3 años.
    \item \textbf{Abonos extraordinarios programados:}
    \begin{itemize}
        \item Final del año 4: 50.000 UVR.
        \item Final del año 5: 50.000 UVR.
    \end{itemize}
\end{itemize}

\textbf{Consideración adicional:} Estimar la inflación durante la vigencia del crédito, dado que la UVR está indexada a esta variable.

\end{frame}

%------------------------------------------------

\begin{frame}{Descripcion del Ejercicio}
\small
El capital inicial debe distribuirse entre los siguientes activos para generar los recursos necesarios para cubrir las obligaciones del crédito.

\begin{center}
\renewcommand{\arraystretch}{1.25}
\begin{tabular}{|p{3.7cm}|p{3.9cm}|p{3.5cm}|}
\hline
\centering\textbf{Activo} & 
\centering\textbf{Rentabilidad anual esperada} & 
\centering\textbf{Liquidez} \tabularnewline
\hline
Bono corporativo & 
\centering 7\% -- 10\% & 
Baja (vence en 3 años) \tabularnewline
\hline
Fondo de inversión conservador & 
\centering 4\% -- 6\% & 
Media (rescate en 30 días) \tabularnewline
\hline
ETF de renta variable & 
\centering 8\% -- 15\% & 
Alta \tabularnewline
\hline
\end{tabular}
\end{center}

\textbf{Problema:} Determinar una estrategia de inversión que permita administrar el capital inicial de \$760.000.000 COP, considerando la rentabilidad, la liquidez y las obligaciones del crédito durante los 10 años.
\end{frame}


%------------------------------------------------
\subsection{Metodologia}
%------------------------------------------------

%------------------------------------------------
\begin{frame}{Metodología — Secuencia de solución}
\small
El ejercicio se resuelve de forma \textbf{secuencial}:

\begin{enumerate}
    \item \textbf{Plan de amortización en UVR:} calcular la cuota fija en UVR,
          aplicar el periodo de gracia y los abonos extraordinarios.
    \item \textbf{Proyección de la UVR} con \textbf{MBG} y \textbf{SARIMAX}
          para convertir cada cuota/abono a pesos \(\rightarrow\) serie \(L_t\).
    \item \textbf{Simulación de activos} (Bono, Fondo, ETF) y estimación de \(\mu\), \(\Sigma\).
    \item \textbf{Optimización Markowitz multiperiodo} sujeta a riesgo y liquidez.
    \item \textbf{Análisis de shortfall} sobre trayectorias Monte Carlo,
          contrastando los flujos \(L_t\) de MBG y SARIMAX.
\end{enumerate}

\vspace{0.2cm}
\begin{center}
\textit{Sin el plan de pagos en UVR no existe el pasivo \(L_t\) que el portafolio debe cubrir.}
\end{center}
\end{frame}


%------------------------------------------------
\begin{frame}{Paso 1 — Plan de amortización en UVR}
\small
Condiciones del crédito y cálculo de la cuota fija en UVR:

\begin{itemize}
    \item Capital: \(C = 1\,000\,000\) UVR; plazo \(n=120\) meses.
    \item Tasa: UVR \(+\,10\%\) E.A. \(\Rightarrow\)
          \(i_m = (1{,}10)^{1/12}-1 \approx 0{,}7974\%\) E.M.
    \item Periodo de gracia: 36 meses con \textbf{cuota \(= 0\)} (el interés se capitaliza).
    \item Abonos extraordinarios: \(50\,000\) UVR en los periodos 48 y 60.
\end{itemize}

\vspace{0.1cm}
La cuota fija sobre el saldo al salir de la gracia (y descontando el valor presente
de los abonos) resulta:

\[
A \approx 20\,379{,}53\ \text{UVR}
\]

En los meses 48 y 60 la salida en UVR es \(A + 50\,000 \approx 70\,379{,}53\) UVR.
Este perfil de pagos en UVR es \textbf{determinístico}; la incertidumbre entra al
convertirlo a pesos con la trayectoria de la UVR.
\end{frame}


%------------------------------------------------
\begin{frame}{Paso 2 — De cuotas UVR a flujo en pesos \(L_t\)}
\small
Cada cuota (o abono) en UVR se convierte a COP con la UVR proyectada del mes \(t\):

\[
L_t = \underbrace{Q_t^{\text{UVR}}}_{\text{plan de amortización}}
      \times \underbrace{\mathrm{UVR}_t}_{\text{MBG o SARIMAX}}
\]

\begin{itemize}
    \item \textbf{SARIMAX:} captura dinámica de corto plazo y exógenas (IPC/TRM).
    \item \textbf{MBG:} proyección estocástica de largo plazo de la UVR.
\end{itemize}

\vspace{0.15cm}
Resultado: dos series de pasivo en pesos,
\(L_t^{\text{MBG}}\) y \(L_t^{\text{SARIMAX}}\),
que alimentan el ALM y el análisis de shortfall.

\end{frame}


%------------------------------------------------
\begin{frame}{Paso 3 — Modelos de proyección de la UVR}
\small
Para obtener \(\mathrm{UVR}_t\) en (2) se usan dos modelos:

\vspace{0.1cm}
\textbf{SARIMAX}
\begin{itemize}
    \item Estacionariedad, órdenes \((p,d,q)\times(P,D,Q)_s\).
    \item Exógenas IPC/TRM (indexación inflacionaria).
\end{itemize}

\textbf{MBG}
\begin{equation*}
S_{t+\Delta t}=S_t\exp\Bigl(\bigl(\mu-\tfrac12\sigma^2\bigr)\Delta t+\sigma\sqrt{\Delta t}\,Z\Bigr)
\end{equation*}
\begin{itemize}
    \item \(\mu\), \(\sigma\) estimados en la serie histórica de la UVR.
\end{itemize}

Ambos generan las trayectorias que, multiplicadas por \(Q_t^{\mathrm{UVR}}\), producen \(L_t^{\mathrm{MBG}}\) y \(L_t^{\mathrm{SARIMAX}}\).
\end{frame}

%------------------------------------------------
\begin{frame}{Paso 4 — Simulación de los activos financieros}
\small
Para cada activo se generan \(K\) trayectorias diarias de 10 años con un proceso de Wiener
escalado y \textbf{acotado} dentro de los rangos de rentabilidad:

\begin{center}
\renewcommand{\arraystretch}{1.15}
\begin{tabular}{|l|c|c|}
\hline
\textbf{Activo} & \textbf{Rango anual} & \textbf{Liquidez} \\
\hline
Bono corporativo & \(7\%{-}10\%\) & Baja (vence en 3 años) \\
Fondo conservador & \(4\%{-}6\%\) & Media (rescate 30 días) \\
ETF renta variable & \(8\%{-}15\%\) & Alta \\
\hline
\end{tabular}
\end{center}

\vspace{0.15cm}
\textbf{Procedimiento:}
\begin{enumerate}
    \item Simulación diaria (Wiener + clipping al rango).
    \item Agregación a retornos mensuales compuestos: \(r_m=\prod(1+r_d)-1\).
    \item Estimación de parámetros de Markowitz:
    \[
    \mu=\mathbb{E}[r_m],\qquad \Sigma=\mathrm{Cov}(r_m).
    \]
\end{enumerate}
\end{frame}


%------------------------------------------------
\begin{frame}{Paso 5 — Portafolio Markowitz y análisis de shortfall}
\small
\textbf{Optimización multiperiodo} (escala mensual):
\begin{equation*}
\begin{aligned}
\max_{w}\quad & w^\top\mu \\
\text{s.a.}\quad
& w^\top\Sigma\,w \le \sigma_{\max}^2,\;
\sum_i w_i = 1,\; w_i\ge 0,\;
w_{\text{Fondo}}+w_{\text{ETF}}\ge\beta.
\end{aligned}
\end{equation*}

\begin{itemize}
    \item P1--P3 (meses 1--108): \{Bono, Fondo, ETF\}
    \item P4 (meses 109--120): \{Fondo, ETF\} \quad\textit{(el bono ya venció)}
\end{itemize}

Dinámica del valor:
\[
V_t = V_{t-1}\,(1+w^\top\mu) - L_t
\]

\vspace{0.1cm}
\textbf{Evaluación de riesgo:} se propagan \(K\) trayectorias de retornos con los pesos óptimos
y se calculan probabilidad de shortfall, Expected Shortfall y la distribución de \(V_{120}\),
para los flujos \textbf{MBG} y \textbf{SARIMAX}.
\end{frame}

%------------------------------------------------
\subsection{Resultados}
%------------------------------------------------

%------------------------------------------------
\begin{frame}{Resultados — Plan de pagos en UVR}
\small
Hitos del plan de amortización (unidad: UVR):

\begin{center}
\renewcommand{\arraystretch}{1.2}
\begin{tabular}{lrr}
\hline
\textbf{Momento} & \textbf{Cuota (UVR)} & \textbf{Saldo final (UVR)} \\
\hline
Meses 1--36 (gracia)     & \(0\)           & \(\approx 1\,331\,000\) \\
Mes 37 (inicia cuota)    & \(20\,379{,}53\)& \(\approx 1\,321\,234\) \\
Mes 48 (+ abono 50\,000) & \(70\,379{,}53\)& \(\approx 1\,158\,530\) \\
Mes 60 (+ abono 50\,000) & \(70\,379{,}53\)& \(\approx 968\,812\) \\
Mes 120 (última cuota)   & \(20\,379{,}53\)& \(0\) \\
\hline
\end{tabular}
\end{center}

\vspace{0.2cm}
\begin{itemize}
    \item Durante la gracia el saldo crece por capitalización de intereses.
    \item La cuota fija \(A \approx 20\,379{,}53\) UVR amortiza el crédito en los 84 meses restantes,
          incorporando los dos abonos extraordinarios.
\end{itemize}
\end{frame}


%------------------------------------------------
\begin{frame}{Resultados — Cuotas en COP según UVR proyectada}
\small
Al multiplicar el plan en UVR por cada trayectoria de UVR se obtiene el pasivo en pesos:

\begin{center}
\renewcommand{\arraystretch}{1.2}
\begin{tabular}{lrr}
\hline
\textbf{Indicador} & \textbf{MBG} & \textbf{SARIMAX} \\
\hline
Cuota típica post-gracia (ej. mes 37) & \(\approx \$9{,}13\) M & \(\approx \$8{,}94\) M \\
Cuota con abono (mes 48)              & \(\approx \$33{,}2\) M & \(\approx \$32{,}3\) M \\
Cuota con abono (mes 60)              & \(\approx \$35{,}1\) M & \(\approx \$33{,}8\) M \\
\textbf{Total pagado (10 años)}       & \(\approx \$986{,}7\) M & \(\approx \$926{,}6\) M \\
\hline
\end{tabular}
\end{center}

\vspace{0.15cm}
SARIMAX proyecta una UVR algo más baja en el horizonte \(\Rightarrow\) pasivo total menor.
Ambas series \(L_t\) son las que enfrenta el portafolio en el ALM y en el shortfall.
\end{frame}

\begin{frame}{Resultados — Parámetros del modelo}
\small
\textbf{Retornos esperados} (estimados a partir de las trayectorias Wiener):

\begin{center}
\begin{tabular}{lcc}
\hline
\textbf{Activo} & \(\mu\) mensual & \(\mu\) anual \\
\hline
Bono corporativo   & \(0{,}672\%\) & \(8{,}37\%\) \\
Fondo conservador  & \(0{,}402\%\) & \(4{,}93\%\) \\
ETF renta variable & \(0{,}896\%\) & \(11{,}29\%\) \\
\hline
\end{tabular}
\end{center}

\vspace{0.2cm}
\textbf{Restricciones:}
\begin{itemize}
    \item \(C_0 = \$760\,000\,000\) COP
    \item \(\sigma_{\text{anual}} \le 8\%\) \(\Rightarrow\) \(\sigma_{\text{mensual}} \le 2{,}31\%\)
    \item Liquidez mínima: \(w_{\text{Fondo}}+w_{\text{ETF}}\ge 30\%\)
    \item \(w_i\ge 0\), \(\sum w_i=1\)
\end{itemize}
\end{frame}


%------------------------------------------------
\begin{frame}{Resultados — Composición óptima del portafolio}
\centering
\includegraphics[width=0.92\textwidth]{images/alm_pesos_por_periodo.png}

\vspace{0.15cm}
\small
Dada la baja volatilidad simulada de los activos, el tope de riesgo no es vinculante y la solución
se concentra en el activo de mayor retorno esperado (ETF), respetando la restricción de liquidez.
\end{frame}


%------------------------------------------------
\begin{frame}{Resultados — Evolución determinista (MBG vs SARIMAX)}
\centering
\includegraphics[width=0.95\textwidth]{images/alm_evolucion_determinista_mbg_sarimax.png}

\vspace{0.1cm}
\small
\begin{center}
\begin{tabular}{lrr}
\hline
 & \textbf{MBG} & \textbf{SARIMAX} \\
\hline
Total pagado del pasivo & \$986{,}7 M & \$926{,}6 M \\
Residual (mes 120)      & \$771{,}4 M & \$851{,}5 M \\
Estado                  & Solvente   & Solvente \\
\hline
\end{tabular}
\end{center}
\end{frame}


%------------------------------------------------
\begin{frame}{Resultados — Robustez Monte Carlo (flujo MBG)}
\centering
\includegraphics[width=0.95\textwidth]{images/alm_percentiles_mbg.png}

\vspace{0.1cm}
\small
\begin{itemize}
    \item \(K\) trayectorias; banda p5--p95 siempre por encima de cero.
    \item Probabilidad de shortfall: \(\mathbf{0\%}\).
    \item Residual simulado (aprox.): media \(\approx\$771\) M; p5 \(\approx\$732\) M.
\end{itemize}
\end{frame}


%------------------------------------------------
\begin{frame}{Resultados — Robustez Monte Carlo (flujo SARIMAX)}
\centering
\includegraphics[width=0.95\textwidth]{images/alm_percentiles_sarimax.png}

\vspace{0.1cm}
\small
\begin{itemize}
    \item El pasivo proyectado por SARIMAX es menor \(\Rightarrow\) residual más alto.
    \item Probabilidad de shortfall: \(\mathbf{0\%}\).
    \item Residual simulado (aprox.): media \(\approx\$852\) M; p5 \(\approx\$811\) M.
\end{itemize}
\end{frame}


%------------------------------------------------
\begin{frame}{Resultados — Distribución del capital residual final}
\centering
\includegraphics[width=0.95\textwidth]{images/alm_hist_residual_mbg_sarimax.png}

\vspace{0.1cm}
\small
En ambos pronósticos de UVR la masa de la distribución de \(V_{120}\) está claramente
por encima de cero: el plan de inversión es \textbf{robusto} ante la incertidumbre de mercado.
\end{frame}


%------------------------------------------------
\begin{frame}{Resultados — Comparación MBG vs SARIMAX}
\centering
\includegraphics[width=0.85\textwidth]{images/alm_comparacion_residuales.png}

\vspace{0.15cm}
\small
\begin{center}
\begin{tabular}{lcc}
\hline
\textbf{Métrica} & \textbf{MBG} & \textbf{SARIMAX} \\
\hline
Prob. shortfall     & \(0\%\) & \(0\%\) \\
Residual p5         & \(\approx\$732\) M & \(\approx\$811\) M \\
Residual mediana    & \(\approx\$771\) M & \(\approx\$851\) M \\
\hline
\end{tabular}
\end{center}
\end{frame}

%------------------------------------------------
\subsection{Discusion}
%------------------------------------------------

\begin{frame}{Discusión}
\small
\begin{enumerate}
    \item \textbf{Plan de pagos:} la cuota cuota fija en UVR
    (\(\approx 20\,379{,}53\)), con gracia sin pago y abonos en 48 y 60,
    define un pasivo determinístico en UVR. La incertidumbre entra solo al
    pasarlo a pesos con MBG y SARIMAX.

    \item \textbf{Pasivo en COP:} SARIMAX implica un total pagado menor
    (\(\approx \$926{,}6\) M) que MBG (\(\approx \$986{,}7\) M). En ambos casos
    el capital inicial de \$760 M, invertido de forma óptima, cubre el flujo
    y deja residual positivo.

    \item \textbf{Portafolio:} bajo la \(\Sigma\) simulada el tope de riesgo
    apenas aprieta y la solución se concentra en el ETF. En una calibración
    con volatilidades de mercado el bono y el fondo tendrían más peso;
    la restricción de liquidez y el vencimiento del bono a 3 años siguen siendo
    las anclas de negocio correctas.

    \item \textbf{Shortfall:} probabilidad \(0\%\) en MBG y en SARIMAX;
    el percentil 5 de \(V_{120}\) permanece holgadamente positivo.
    El plan es robusto ante la incertidumbre de la UVR y de los retornos.
\end{enumerate}
\end{frame}


\begin{frame}{Conclusión del ejercicio ALM}
\small
\begin{center}
\fbox{\parbox{0.92\textwidth}{\centering
Partiendo del plan de amortización en UVR, convertido a pesos con MBG y SARIMAX,
el portafolio Markowitz multiperiodo cubre íntegramente las cuotas y los abonos
extraordinarios a 10 años, con residual esperado positivo y sin evidencia de
shortfall en las simulaciones realizadas.
}}
\end{center}

\vspace{0.35cm}
\textbf{Aportes del ejercicio}
\begin{itemize}
    \item Integración pasivo (UVR) -- activo (Markowitz) en un mismo marco ALM.
    \item Uso de dos proyecciones de UVR para contrastar sensibilidad del plan.
    \item Evaluación no solo en media, sino en distribución (percentiles y shortfall).
\end{itemize}
\end{frame}


\section{Referencias}
\begin{frame}{Referencias}
\begin{thebibliography}{99}

\bibitem{barriga2021}
Barriga Manrique, E. (2021).
\newblock \textit{Matemáticas financieras}.
\newblock Notas de clase no publicadas, Departamento de Finanzas, Universidad EAFIT.

\end{thebibliography}
\end{frame}

\end{document}